\section{Current dependency audit: the six Hida fields and both support installations}
\label{sec:dependency-audit-v45}

\subsection{Current status and reading order}

Version 0.45 completes the six fields of the original
\texttt{HidaPropsAlt} on one ideal $J$, constructs the actual
all-level \texttt{LevelAlt} and its \texttt{HasAut}, and installs
the selected rank-four Hida/$P_1$ support result on both original
Links above one common coefficient threshold.  The entire selected
\texttt{COUNT\_II\_SUPPORT} obligation is consequently discharged.
Its original field, embedding, residual, determinant and
pro-$\ell$ data are retained, including $\ell=3$.

The pinned Alpoge--Claude formal proof is accepted as fully checked
and correct.  This remains a \textbf{partial human reconstruction}
of RH and unrestricted Dirichlet GRH.  The common geometric and
arithmetic families, global estimates and several case producers
still have open human reconstruction contracts.  The all-rank
\texttt{HidaUAlt} conclusion is distinguished from full
\texttt{ModUAlt}: the latter still requires the original
\texttt{PfromAut} supplier, including arbitrary-chain $P_3$.
No Lean build or kernel replay is claimed.

There are @NEW@ new and @TOTAL@ cumulative claims.  All @OLD@
predecessor claim records, source arrays and dependency edges are
preserved.  The selected sufficient-route graph still has 121 nodes
and 161 edges.  The whole-obligation register decreases from 31 to
30: 21 class-C and 9 class-B tasks.  Of these, 23 are shared by RH
and Dirichlet GRH, and 7 are additional for Dirichlet GRH.  The
claim inventory is not a percentage of proof completion.

HED and HAW prove arithmetic-point decomposition and automorphy
away from the exceptional set.  HPT returns an arbitrary supplied
automorphic point to the actual Hecke support; HTS proves twisting
on that same support.  HLM constructs the complete family of
auxiliary-level modules, and HTW constructs its integral
Taylor--Wiles data.  HIA assembles all six Hida fields and performs
the two selected support installations.  The separate reader
contains this audit and these seven chapters.  The cumulative
edition preserves the previous proof body verbatim.  CMP1--CMP6
remain the complete conditional terminal deduction on their
original hypotheses.

\subsection{The same ideal and all six Hida properties}

The empty auxiliary level gives an actual ordinary Hecke module
and its original \texttt{IsJ} ideal.  Every subsequent field refers
to that ideal, the same coefficient ring and the same universal
determinant.  The nonzero module branch uses the kernel of the
actual Hecke-algebra map; the zero branch retains the top ideal.

\begin{center}
\begin{tabular}{@{}p{0.13\textwidth}p{0.53\textwidth}p{0.22\textwidth}@{}}
\toprule
Original field & Human-readable output & Completed supplier\\
\midrule
\texttt{ring} & All original quotient-ring and determinant properties & JRP15\\
\texttt{dense} & Bare arithmetic-point density for this same $J$ & UPA18\\
\texttt{pt} & The original point theorem, with its rank-$2$--$4$ antecedent & HPT18\\
\texttt{end1} & Full arithmetic-point decomposition & HED16\\
\texttt{end2} & Full irreducible arithmetic-point automorphy away & HAW19\\
\texttt{twist} & The original all-positive-rank twisting property & HTS8\\
\bottomrule
\end{tabular}
\end{center}

HED starts with the given arithmetic point and its given continuous
semisimple representation.  Surjective quotient factorization
preserves the point, weight and full polynomial law after every
scalar extension.  The base-change decomposition supplies one
initial Galois representation for every cusp and every positive
Speh length.  Exact inverse cyclotomic strings retain the
Frobenius normalization.  Degree-one Chebotarev turns the cofinite
Frobenius product into an identity of characteristic polynomials
at every group element.

When all string lengths are one, the already matched
generic-unitary classification gives the strict spherical interval
\[
                   q^{-1/2}<|\alpha|<q^{1/2}.
\]
It yields the original cofinite strict ratio
$|z|<q|z'|$ for every supplied Frobenius multiset at every required
place.  This branch does not discard the other positive lengths
from the full decomposition theorem.

HAW proves the original ordinary principal-series description at
every arithmetic point at its literal scope.  A nonzero Jacquet
character block has a character quotient by finite-length
Nakayama; a functional need not detect a preselected eigenvector.
The ordinary projector selects coefficient-field units, including
the exclusion of every nonunit.  The ordered flag and conjugate
place comparison retain the polarization multiplier
$\epsilon^{1-n}$ and the reversal of dual characters.

One global real norm twist unitarizes the central character before
any finite place is selected.  It therefore gives one common center
for all the required local exponent windows.  The exact local
character calculation then shows that a forbidden positive-power
cyclotomic ratio would force a factor $q$ between two character
norms.  The strict window excludes this.  The resulting
\texttt{AutOffN} keeps the source's conditional continuity in
its original location.

HPT begins with the supplied automorphic representation and
constructs a nonzero vector in the actual ordinary space, with
the integral type, infinity weight, depth and residual block all
retained.  Its Hecke point equals the supplied point because the
full determinant and ordered characters agree in a finite
coefficient ring, where the universal-ring uniqueness theorem
applies.  Thus $J\subseteq\ker x$ is the conclusion.  The
all-positive-rank return theorem is separate from the last point
field, whose original antecedent includes $2\le n\le4$ and
whose use of \texttt{DescWeak} remains explicit.

HTS realizes a permitted odd-order conjugate-dual twist by an
actual adelic unitary determinant character.  If its order is
$m=2N+1$, the power $\mu^{N+1}$ is a square root because
$2(N+1)\equiv1\pmod m$.  At a split place the resulting
determinant character has precisely the required square.
Multiplication of the actual automorphic functions preserves the
infinity weight and the required local vectors.  The guarded
local comparisons and full polynomial-law twist supply every
input of the all-rank point-return theorem.  This proves
\texttt{TwistOk} without importing the low-rank descent premise.

\subsection{One family over every auxiliary level}

HLM constructs the actual directed union of smooth automorphic
modules and identifies its fixed modules at every permitted
level.  Integral trace projectivity supplies a group-algebra
splitting without dividing by the quotient-group order.  The
ordinary and residual projectors preserve the required module
properties, with the residual decomposition justified by
Artin--Rees.

At every level the determinant law is constructed before using
any point theorem which assumes it.  Compact interpolation is
performed over the coefficient ring; integral word identities
are tested in finite-degree truncated scalar algebras.  Polynomial
coefficient separation establishes the complete law after
arbitrary scalar extension.  The full ordinary pair then gives
the surjective universal-ring map.  Universal uniqueness glues
these maps into one coherent action on the directed union.

Classical points separate actual module vectors.  The type and
local correspondence comparisons retain all of the original
quantifiers.  Dense Frobenius comparison identifies the old-level
residual component with the initial Hecke module, preserving
$J$.  One base is chosen before all auxiliary level data,
including the empty auxiliary set.  This proves the literal
\texttt{LevelAlt} in every positive rank.

\subsection{Integral Taylor--Wiles comparison and completion}

HTW retains the diagonal-unit terms in the tame Hecke relation.
For a simple reflection it has the form
\[
 T_s^2=Q+T_s\sum_c D_c,\qquad T_sD_c=D_c^sT_s.
\]
The corresponding left and right inverses have their respective
unit terms in the correct order.  The usual Iwahori relation is
the specialization in which those diagonal terms act trivially.
The resulting torus operators and an explicit Jacquet dimension
bound give one exponent $n!$, independent of the later auxiliary
choices and of isotypic multiplicity.

The comparison of spherical and Iwahori modules is integral.
For ordered variables put
\[
 d=\prod_{i<k}(X_k-X_i),\qquad
 N_Q=\prod_{i<k}(X_k-QX_i).
\]
An explicit finite family satisfies
\[
              \sum_T a_T\,w(b_T)=\delta_{w,1}\,d.
\]
With the divided difference $\partial=\mathrm{Alt}/d$, these
identities give both dual expansions over the symmetric ring.
The actual Hecke trace operator is
$E(f)=\partial(N_Qf)$; its integral identity is proved through
the faithful polynomial representation.  These formulas construct
both comparison maps.  Multiplying by the remaining prescribed
factors makes both compositions exactly the original ordered-pair
\texttt{fEl}.  No inversion of $n!$ or replacement by an
unspecified annihilator is made.

The desired topologically nilpotent part is a canonical summand
over the finite coefficient algebra.  Its projectivity makes it
finite free over the local diamond group algebra.  Power actions
extend coherently through the inverse limit.  Reduction by the
diamond ideal gives the correct ordered-root summand, and
completion preserves the two comparison maps.  The zero-module
case is handled separately when using the finite image algebra.

The coefficient map is also explicit: the root variables map
to the selected roots minus their prechosen residue lifts, and
the additional variables map to zero.  Its composition with the
coefficient embedding equals the original prescribed map, not
merely a map agreeing on closed points.  Every field of
\texttt{TWRing} and \texttt{TWMaps}, both full trace relations,
and \texttt{HasAut} is obtained.  One base and one lift of each
residual root are fixed before all later $N$.

\subsection{Both rank-four support installations}

HIA uses the empty-level \texttt{IsJ} once, attaches the six
fields listed above to it, and attaches the same $J$'s
\texttt{HasAut}.  The inherited selected patching, normalization,
pro-$\ell$ and $P_1$ consumers then yield the required rank-four
support condition.  Their original constants and choice order
are preserved; no arbitrary-choice or arbitrary-chain
strengthening is used.

The original two QX Links are installed with fixed $\iota$ above
one common coefficient threshold, before the later coefficient
DVR, universal ring and weight algebra are quantified.  Full
\texttt{CohP}, \texttt{LocCFT}, the original residual
cyclotomic condition, determinant relation and pro-$\ell$
conditions remain visible.  This is the precise selected closure
recorded by \texttt{COUNT\_II\_SUPPORT}.

The general six-field \texttt{HidaUAlt} is also obtained for
all positive ranks: outside ranks two through four the point
field has an impossible antecedent.  Full \texttt{ModularAlt}
and \texttt{ModUAlt} retain the original \texttt{PfromAut}
premise.  The selected rank-four $P_1$ conclusion is not used
as a substitute for arbitrary-chain $P_3$.

\subsection{Every remaining selected obligation}

Each contract below is an entire selected output.  The other 30
contracts retain their original scope and dependency edges.
The next principal work is the prescribed common geometric and
arithmetic family, its three sufficient cylinder counts, the
global quantitative estimates on that family, and the remaining
one-character, congruence, rigidity and case producers.  The named
base-change, attachment, local, descent, cohomological and period
foundations remain explicit at their recorded scopes.  Current
frontier descriptions supersede historical notes in the preserved
chapters.
